% GENERATED FILE. Edit canonical lessons, then run npm run generate:books.
% canonical-source-hash: fnv1a64:a89eb764aebea74c
% canonical-lessons: ES-C410-academia, ES-C410-frances, ES-C410-examen, ES-C410-education-recall-1, ES-C410-education-recall-2

\chapter{Academy, French, Student and Exam}
\label{ch:academia-frances-estudiante-examen}
\hlchaptermodality{410}

\begin{chapteropening}
\textbf{By the end of this chapter:} \emph{I can identify an academy, French, a student and an exam in simple education messages.}

Retrieve four education words across course, identity and assessment messages.
\end{chapteropening}

\section[la academia]{la academia — from a place-name to a place of learning}

\label{lesson:ES-C410-academia}



\begin{quote}
Recall \textbf{estudiante}, a student. A student may learn at \textbf{la academia}, an academy or teaching centre.
\end{quote}

\begin{cousinweb}
Spanish \textbf{academia} came through Latin \textbf{academia} from Greek \textbf{Akadēmía}, the place-name associated with Plato's school. English \textbf{academy} belongs to the same travelled family. In an A1 message, \textbf{una academia de idiomas} is a language school, not necessarily a university institution.
\end{cousinweb}

\begin{grammarlens}[title={Grammar Lens: what kind of academy?}]
Use \textbf{de} to name the subject: \textbf{una academia de idiomas}. The noun is feminine: \textbf{la academia}, \textbf{una academia}.
\end{grammarlens}

\subsection*{Your turn}
Say these aloud:

\begin{itemize}
\raggedright
  \item \emph{la academia}
  \item \emph{una academia de idiomas}
  \item \emph{soy estudiante en una academia}
\end{itemize}

\begin{tcolorbox}[breakable,colback=teal!4,colframe=teal!35!black,title={Before you move on}]
Academy or teaching centre? (\emph{La academia}.)

Source: \href{https://dle.rae.es/academia}{DLE: academia}.
\end{tcolorbox}
\section[francés / francesa]{francés / francesa — language and person}

\label{lesson:ES-C410-frances}



\begin{quote}
At \textbf{una academia de idiomas}, one possible subject is \textbf{francés}, French.
\end{quote}

\begin{cousinweb}
Spanish \textbf{francés} came from French \textbf{français}. The language name is the masculine singular form: \textbf{hablo francés}. For people and adjectives, the ending agrees: \textbf{un estudiante francés}, \textbf{una estudiante francesa}.
\end{cousinweb}

\begin{grammarlens}[title={Grammar Lens: the accent comes and goes}]
The final stress of \textbf{francés} needs its written accent. Adding \textbf{-a} moves the natural stress: \textbf{francesa}, with no accent mark. The plural follows the same pattern: \textbf{franceses, francesas}.
\end{grammarlens}

\subsection*{Your turn}
Say these aloud:

\begin{itemize}
\raggedright
  \item \emph{francés}
  \item \emph{una academia de francés}
  \item \emph{una estudiante francesa}
\end{itemize}

\begin{tcolorbox}[breakable,colback=teal!4,colframe=teal!35!black,title={Before you move on}]
French? (\emph{Francés}.) A French woman or feminine adjective? (\emph{Francesa}.)

Source: \href{https://dle.rae.es/franc\%C3\%A9s}{DLE: francés}.
\end{tcolorbox}
\section[el examen]{el examen — a weighing of what you know}

\label{lesson:ES-C410-examen}



\begin{quote}
A student of French may take \textbf{un examen de francés}, a French exam.
\end{quote}

\begin{cousinweb}
Spanish \textbf{examen} is Latin \textbf{exāmen}, a weighing or examination. English \textbf{exam} and \textbf{examination} are learned relatives. The old weighing image is a useful memory hook: an exam weighs what has been learned.
\end{cousinweb}

\begin{grammarlens}[title={Grammar Lens: singular and plural stress}]
The singular is \textbf{el examen}, with no written accent because the stress falls naturally on \textbf{-xa-}. The plural is \textbf{los exámenes}; the accent appears to keep that stress when another syllable is added.
\end{grammarlens}

\subsection*{Your turn}
Say these aloud:

\begin{itemize}
\raggedright
  \item \emph{el examen}
  \item \emph{un examen de francés}
  \item \emph{los exámenes}
\end{itemize}

\begin{tcolorbox}[breakable,colback=teal!4,colframe=teal!35!black,title={Before you move on}]
Exam? (\emph{El examen}.) Exams? (\emph{Los exámenes}.)

Source: \href{https://dle.rae.es/examen}{DLE: examen}.
\end{tcolorbox}
\section[(review)]{education recall one — place, subject, learner, assessment}

\label{lesson:ES-C410-education-recall-1}



\begin{quote}
Give the place, subject, learner and assessment. (\emph{Academia, francés, estudiante, examen}.)
\end{quote}

\subsection*{Your turn}
Say these aloud:

\begin{itemize}
\raggedright
  \item \emph{una academia de idiomas}
  \item \emph{hablo francés}
  \item \emph{soy estudiante}
  \item \emph{un examen de francés}
\end{itemize}

\begin{tcolorbox}[breakable,colback=teal!4,colframe=teal!35!black,title={Before you move on}]
Academy, French, student, exam? (\emph{Academia, francés, estudiante, examen}.)
\end{tcolorbox}
\section[(review)]{education recall two — complete the message}

\label{lesson:ES-C410-education-recall-2}



\begin{quote}
Complete each education message from meaning.
\end{quote}

\subsection*{Your turn}
\begin{itemize}
\raggedright
  \item “It is a language academy.” {[}YOU SAY: \emph{Es una academia de idiomas.}{]}
  \item “I speak French.” {[}YOU SAY: \emph{Hablo francés.}{]}
  \item “I am a student.” {[}YOU SAY: \emph{Soy estudiante.}{]}
  \item “The exam is difficult.” {[}YOU SAY: \emph{El examen es difícil.}{]}
\end{itemize}

\begin{tcolorbox}[breakable,colback=teal!4,colframe=teal!35!black,title={Before you move on}]
Say the four education words once more. (\emph{Academia, francés, estudiante, examen}.)
\end{tcolorbox}
